\section{Conclusion}
\label{sec:conclusion}

We identify a linearly readable human-AI self-representation that emerges
during pretraining and can be manipulated directly at inference time. Across
OLMo 2 checkpoints, the identity contrast forms within the first 3.7\%
of stage-1 pretraining, before instruction tuning. Post-training later rotates the
direction rather than creating the contrast, while on Gemma 3 4B, instruction
tuning makes AI the default self-representation. Writing register strongly
affects raw projection, but within our bounded Gemma 3 4B specificity check,
register steering does not reproduce the behavioral effect and identity
remains the best-supported interpretation of what axis steering changes.

The representation is also context-sensitive: across three models from two
families, a user's identity claim shifts the model's own self-representation.
Steering changes self-report much more strongly than
behavior, but its behavioral effects can still be safety-relevant.
Human-ward steering increases willingness to cause harm significantly on
Gemma 3 4B and Qwen 2.5 7B, and directionally on 12B, while also raising
prompt-injection execution 16-fold on Gemma 3 4B. These effects vary by
model, task, dose, and steering direction, and do not establish a general safety improvement or
degradation. Together, these results suggest that human-versus-AI
self-representation is a manipulable inference-time state whose behavioral
consequences are real but not captured by self-report alone.
